\section{Predictive and Analytical Extensions}
\label{sec:future:extensions}

Three predictive and analytical extensions follow naturally from the historical foundation of the proposed system. The extensions consume the existing aggregate tables of Layer 4 and require no revisiting of the lower warehouse layers. Each is recorded as out of scope for Phase Two but feasible against the same data the proposed system already uses.

\subsection{Real-Time Arrival-Time Prediction}

Lin and Zeng~\cite{linzeng2014} demonstrate a link-node route representation that projects Global Positioning System (GPS) pings onto a one-dimensional route path, constructs per-trip time-distance trajectories, and predicts arrival time at downstream stops. The dominant predictor in their model is the dwell time at time-check stops. The same construction applies to the present project's GPS data once the full fifty-day window is loaded. The extension lifts the time-distance trajectory diagram of the Route-Level Analysis theme (Section~\ref{sec:solution:viz}) from a backward-looking visualisation to a predictive one, with arrival-time bands rendered on the per-route detail view.

\subsection{Per-Route Anomaly Detection}

The anomaly view in the Network Health and Performance theme operates at the network level using a tunable threshold over fixed quantile bands. The natural extension flags anomalies at the per-route per-hour level using statistical control limits derived from the ninety-day Automated Fare Collection (AFC) history. Control charts at three standard deviations on the aggregate-by-route-hour table are sufficient for the construction. The extension feeds back into the Network Health and Performance theme without requiring a new aggregate table or any change to the warehouse schema.

\subsection{Bus Telemetry as a Traffic-Sensing Substrate}

Derevitskiy and colleagues~\cite{derevitskiy2016} aggregate bus GPS pings into road-segment edge-speed estimates and propagate them to road segments where no buses currently operate, using a macroscopic travel-time model. The reported mean absolute percentage error is around twenty percent at distances under fifteen kilometres. The construction surfaces city-wide road-segment travel times as a derived analytical product. The extension is feasible against the full fifty-day GPS dataset but is treated as out of scope for Phase Two because it introduces an analytical theme that is not part of the catalogue in Chapter~\ref{chapter:solution}; the construction is recorded as a candidate for a follow-on phase if the project continues beyond Phase Two.
